\subsection{Distance Measures}\label{subsec:background_distance}

This subsection defines the various distance measures used in cosmology. These are essential for converting observations (angles on the sky, apparent magnitudes) into physical quantities.

\subsubsection{Comoving Distance}

The comoving distance $\chi(z)$ to an object at redshift $z$ is the distance accounting for the expansion of the universe. It is defined as:

\begin{equation}
\chi(z) = \int_0^z \frac{c \, dz'}{H(z')} = \int_0^z \frac{dz'}{E(z')}
\label{eq:comoving_distance}
\end{equation}

(In natural units where $c = H_0 = 1$.)

At large redshift, $\chi(z)$ approaches $\chi_\infty \approx 2 \, c / H_0 \approx 14,000 \, \mathrm{Mpc}$ (the horizon distance).

\remark{
Comoving distance is the natural distance scale for cosmological surveys because it removes the effect of cosmic expansion. Two galaxies at the same redshift but different angular positions on the sky have comoving separation that is independent of when we observe them.
}

\subsubsection{Angular Diameter Distance}

The angular diameter distance $D_A(z)$ relates the angular size $\theta$ (in radians) to the physical size $s$ of an object:

\begin{equation}
\theta = \frac{s}{D_A(z)}
\end{equation}

For a flat universe:
\begin{equation}
D_A(z) = \frac{\chi(z)}{1+z}
\label{eq:D_A}
\end{equation}

The factor $(1+z)$ accounts for the redshift dimming: sources at higher redshift appear smaller and dimmer.

\subsubsection{Luminosity Distance}

The luminosity distance $D_L(z)$ relates the absolute luminosity $L$ of a source to the observed flux $F$:

\begin{equation}
F = \frac{L}{4\pi D_L(z)^2}
\end{equation}

For a flat universe:
\begin{equation}
D_L(z) = (1+z) \chi(z)
\label{eq:D_L}
\end{equation}

Note that $D_L(z) = (1+z)^2 D_A(z)$. This duality is a consequence of the reciprocity theorem in general relativity.

\subsubsection{Comoving Volume Element}

The comoving volume element per unit solid angle per unit redshift is:

\begin{equation}
\frac{dV_c}{d\Omega \, dz} = \frac{\chi(z)^2}{H(z)}
\label{eq:dVc}
\end{equation}

This is essential for computing the number density of galaxies and other volume statistics.

\example{
For $\Om = 0.3, \OL = 0.7, H_0 = 70 \, \mathrm{km \, s^{-1} \, Mpc^{-1}}$:
\begin{itemize}
    \item $\chi(0.5) \approx 1547 \, \mathrm{Mpc}$
    \item $D_A(0.5) \approx 1031 \, \mathrm{Mpc}$
    \item $\frac{dV_c}{d\Omega \, dz}\big|_{z=0.5} \approx 1.1 \times 10^9 \, \mathrm{Mpc}^3 \, \mathrm{deg}^{-2}$
\end{itemize}
}

\subsubsection{Relating Observations to Physical Quantities}

\begin{itemize}
    \item \textbf{Angular scale on sky}: A small angle $\theta$ corresponds to physical scale $s = \theta D_A(z)$
    \item \textbf{Apparent magnitude}: Related to luminosity distance via $m - M = 5 \log_{10} D_L(z) - 5$
    \item \textbf{Galaxy counts}: Number of galaxies in a redshift bin depends on comoving volume via \cref{eq:dVc}
\end{itemize}

These distance measures are fundamental inputs to weak lensing kernels (\cref{ch:weak_lensing}) and observables computations (\cref{ch:observables}).

See \cref{subsec:background_model} for the definition of $E(z)$ and \cref{subsec:background_evolution} for lookback time.
